\documentclass[11pt, a4paper]{article}
\usepackage[a4paper, top=2.5cm, bottom=2.5cm, left=2cm, right=2cm]{geometry}
\usepackage{amsmath}
\usepackage{amssymb}
\usepackage[colorlinks=true, linkcolor=blue, urlcolor=blue, citecolor=blue]{hyperref}

\title{Omega Theory --- Volume 14\\\Large The Dark Sector: A Model-Scoped Summary}
\author{Omega Protocol}
\date{\today}

\begin{document}

\maketitle

\begin{quote}
\textbf{Status.} Formal counterpart: \texttt{lean\_proofs/Vol14\_DarkSector.lean}. The file is a model. It establishes: (i) a cosmic density budget on the 2-simplex, in which the dark-energy fraction is determined by the matter fractions for any nonnegative fractions summing to one; (ii) a strict bound, under which the dark-energy fraction is below one whenever the baryonic and dark-matter fractions are strictly positive; (iii) a concrete $\Lambda$CDM-type benchmark in which dark energy dominates. The file does not construct dark matter or dark energy from the framework. The ``informational pressure'' and ``gravitational memory'' mechanisms described in earlier text have no formal counterpart here. Physical interpretation: open.
\end{quote}

\section*{Layman's Summary}
Observations indicate that most of the energy content of the universe is not ordinary matter. Dark matter is inferred from the gravitational behavior of galaxies, and dark energy from the accelerated expansion. This volume states the bookkeeping constraints that any such budget must satisfy, and a concrete benchmark. It does not identify what dark matter or dark energy is.

\section{Dark Energy as Budget Remainder (Formal Result within the Model)}
If the baryonic, dark-matter, and dark-energy fractions are nonnegative and sum to one, the dark-energy fraction is fixed by the other two. The formal file also shows that it is below one whenever the other two are strictly positive. This is an algebraic fact about a budget, not a mechanism. The earlier claim that dark energy is the baseline informational pressure of the vacuum is withdrawn; it has no formal counterpart.

\section{Dark Matter Mechanism (No Formal Counterpart)}
The earlier claim that space ``remembers'' past motion and produces an extra gravitational pull mimicking dark matter has no formal counterpart in this volume. It is withdrawn as a claim of this work and is listed as an open item.

\section{Modified Newtonian Dynamics (No Formal Counterpart)}
The earlier claim that entropic bounds reproduce MOND-like behavior has no formal counterpart. The file does not derive galactic rotation curves. This is an open item.

\section{Scope and Open Items}
\begin{itemize}
  \item Not established here: the physical nature of dark matter or dark energy; galactic rotation curves; agreement with observed cosmological parameters beyond the benchmark.
\end{itemize}

\end{document}
